
\subsubsection{3.1 Overview}

The study addresses the following question: which feedback signal --- execution test results or pylint/mypy static analysis warnings --- produces larger pass@1 improvement when both are given the same inference compute budget? The iso-compute constraint is the methodological core of the design. Without it, apparent feedback quality differences may simply reflect differences in the number of tokens spent on repair. The B=1000 output token budget is held constant across all conditions.

\subsubsection{3.2 Experimental Conditions}

Three conditions are evaluated on HumanEval (164 problems) and MBPP (378 problems):

\textbf{Condition A (Baseline): No-feedback single-pass generation.} Code is generated once via greedy decoding. No repair loop is applied. All token budget B is consumed by the initial generation.

\textbf{Condition B (Pylint/mypy): Static analysis iterative repair.} After initial generation, failing solutions receive pylint and mypy output formatted as structured prompt context. The model re-generates given the original problem statement, the failed solution, and the static analysis feedback. Repair continues until the token budget is exhausted or the solution passes.

\textbf{Condition C (Execution): Execution test feedback iterative repair.} After initial generation, failing solutions are executed against the benchmark test suite within a sandboxed subprocess. The error output --- exception type, traceback (truncated to $\leq 512$ characters), and first failing assertion --- is formatted as structured prompt context. The model re-generates given the original problem statement, the failed solution, and the execution feedback.

All three conditions use Llama 3.1 8B Instruct with greedy decoding (temperature=0.0, seed=42) and token budget B=1000 output tokens per problem.

\subsubsection{3.3 Iso-Compute Token Budget}

The token budget B=1000 output tokens per problem serves as the experimental unit of compute. Repair is halted when the remaining budget falls below 50 tokens or when the solution passes. At B=1000, with max\_tokens=512 per generation round, initial generation typically consumes 400--512 tokens, leaving at most one repair round before budget exhaustion. Per-round trajectory data from both conditions confirms that rounds 2 and 3 contribute near-zero incremental improvement. Results therefore characterize the effect of a single repair round of each feedback type.

\subsubsection{3.4 Feedback Signal Construction}

\textbf{Pylint/mypy feedback.} \texttt{pylint --output-format=text} is run with the default configuration (all categories enabled) and \texttt{mypy} is run on each failing solution. Output is parsed and provided as structured prompt context to the repair prompt.

\textbf{Execution feedback.} The failing solution is executed in a sandboxed subprocess with a 15-second timeout against the benchmark's test assertions. The exception type, truncated traceback, and first failing assertion are formatted as structured repair context.

\subsubsection{3.5 Statistical Analysis}

\textbf{Primary statistical test:} McNemar's test on the paired $2\times 2$ contingency table (Condition C passes, Condition B passes), at significance level $\alpha =0.05$. The test uses an exact binomial variant when the number of discordant pairs is fewer than 25 (HumanEval), and the chi-square approximation otherwise (MBPP).

\textbf{Effect size:} $\Delta _pass$@1 = pass@1(condition) $-$ pass@1(no-feedback baseline).

\textbf{Bootstrap confidence intervals:} 95\% intervals, 10,000 samples, seed=42.

\subsubsection{3.6 Pylint Coverage Analysis}

For the 64 HumanEval baseline failures, pylint and mypy are run on each failing solution without execution. All pylint flags are recorded by category (E: Error, W: Warning, C: Convention, R: Refactor, I: Information). Total coverage, functional coverage (E+W categories only), and the category distribution of all flags are computed. Bootstrap 95\% CIs for coverage fractions are computed using 10,000 samples, seed=42.

\subsubsection{3.7 Model and Infrastructure}

\begin{table}[htbp]
\centering
\resizebox{\columnwidth}{!}{%
\begin{tabular}{ll}
\toprule
Parameter & Value \\
\midrule
Model & Llama 3.1 8B Instruct (meta-llama/Llama-3.1-8B-Instruct) \\
Backend & vLLM v0.10.1.1 (bfloat16) \\
Hardware & $5\times$ H100 NVL GPUs \\
GPU memory utilization & 0.4 \\
Max model length & 4096 tokens \\
Token budget B & 1000 output tokens per problem \\
Max repair rounds & 3 (effectively 1 at B=1000) \\
Execution sandbox timeout & 15 seconds per test case \\
Decoding & Greedy (temperature=0.0, seed=42) \\
\bottomrule
\end{tabular}
}
\end{table}

